% Motor: XeLaTeX. Sesión 30 - material docente UTP.
% Estructura docente de la Sesión 29.
\documentclass[aspectratio=169,10pt,c]{beamer}
% Reutiliza sin cambios la identidad visual UTP validada en la Sesion 08.
% Reutiliza sin cambios la identidad visual UTP validada en la Sesion 05.
% La segunda ruta permite compartir los assets comunes sin duplicar binarios.
% !TEX root = slides.tex
% Preambulo comun: estilo UTP, colores, fuentes, plantillas de beamer y
% comandos de apoyo didactico. Se carga desde slides.tex tras \documentclass.

% La opción "c" de beamer reparte el espacio libre 1:1.5 (arriba:abajo), por lo
% que el contenido queda por encima del centro. Igualamos ambos rellenos.
\makeatletter
\define@key{beamerframe}{c}[true]{%
  \beamer@frametopskip=0pt plus 1fill\relax%
  \beamer@framebottomskip=0pt plus 1fill\relax%
  \beamer@frametopskipautobreak=0pt plus .5\paperheight\relax%
  \beamer@framebottomskipautobreak=0pt plus .5\paperheight\relax%
}
\makeatother

\usepackage{graphicx}
\usepackage{tikz}
\usepackage{xcolor}
\usepackage{iftex}

\ifPDFTeX
  \usepackage[utf8]{inputenc}
  \usepackage[T1]{fontenc}
  \usepackage{lmodern}          % tipografia escalable para el respaldo con pdflatex
  \renewcommand{\familydefault}{\sfdefault}
\else
  \usepackage{fontspec}
  \IfFontExistsTF{Arial}
    {\setsansfont{Arial}[BoldFont={Arial Bold}]}
    {\setsansfont{Helvetica}[BoldFont={Helvetica Bold}]}
\fi

\usepackage[spanish,es-nodecimaldot]{babel}
\usepackage{microtype}
\usepackage{booktabs}
\usepackage{array}
\usetikzlibrary{arrows.meta,positioning,fit,calc,shapes.geometric}

\graphicspath{{assets/}}

% ---------- Identidad visual UTP ----------
\definecolor{utpRojo}{HTML}{C42B35}
\definecolor{utpNegro}{HTML}{111111}
\definecolor{utpGris}{HTML}{5F5F5F}
\definecolor{utpGrisMedio}{HTML}{9A9A9A}
\definecolor{utpGrisClaro}{HTML}{F1F1F1}

\usetheme{default}
\setbeamercolor{normal text}{fg=utpNegro,bg=white}
\setbeamercolor{structure}{fg=utpRojo}
\setbeamercolor{frametitle}{fg=utpNegro,bg=white}
\setbeamercolor{title}{fg=utpNegro}
\setbeamercolor{subtitle}{fg=utpNegro}
\setbeamercolor{author}{fg=utpNegro}
\setbeamercolor{date}{fg=utpNegro}
\setbeamercolor{block title}{fg=white,bg=utpNegro}
\setbeamercolor{block body}{fg=utpNegro,bg=utpGrisClaro}
\setbeamercolor{alertblock title}{fg=white,bg=utpRojo}
\setbeamercolor{alertblock body}{fg=utpNegro,bg=utpRojo!7}
\setbeamercolor{exampleblock title}{fg=white,bg=utpGris}
\setbeamercolor{exampleblock body}{fg=utpNegro,bg=utpGris!8}
\setbeamercolor{terminal}{fg=white,bg=utpNegro}

\setbeamerfont{title}{series=\bfseries,size=\fontsize{19.5}{22.5}\selectfont}
\setbeamerfont{subtitle}{size=\fontsize{11.5}{13.5}\selectfont}
\setbeamerfont{date}{series=\bfseries,size=\fontsize{11.5}{13.5}\selectfont}
\setbeamerfont{author}{size=\fontsize{10.5}{12.5}\selectfont}
\setbeamerfont{frametitle}{series=\bfseries,size=\large}
\setbeamerfont{block title}{series=\bfseries,size=\small}
\setbeamerfont{block body}{size=\normalsize}

\setbeamertemplate{navigation symbols}{}
\setbeamertemplate{headline}{}
\setbeamertemplate{blocks}[default]
\setbeamertemplate{itemize item}{\color{utpRojo}\small\raise.15ex\hbox{$\blacktriangleright$}}
\setbeamertemplate{itemize subitem}{\color{utpGris}\scriptsize$\bullet$}
\setbeamertemplate{enumerate item}{\color{utpRojo}\insertenumlabel.}
\setbeamertemplate{background canvas}{\color{white}\rule{\paperwidth}{\paperheight}}

% Título de diapositiva: negrita negra a la izquierda, logo UTP a la derecha
\setbeamertemplate{frametitle}{%
  \begin{tikzpicture}[remember picture,overlay]
    \node[anchor=north east,inner sep=0pt]
      at ([xshift=-7mm,yshift=-5.5mm]current page.north east)
      {\includegraphics[width=17mm]{utp-logo.png}};
  \end{tikzpicture}%
  \vspace{1.2mm}%
  \begin{minipage}[t]{.80\textwidth}\raggedright\insertframetitle\end{minipage}\par
  \vspace{1.1mm}{\color{utpRojo}\rule{\textwidth}{1.1pt}}
}

% Solo numeración discreta; la plantilla UTP no lleva pie de página
\setbeamertemplate{footline}{%
  \hfill{\color{utpGrisMedio}\scriptsize\insertframenumber}\hspace{6mm}\vspace{2.5mm}%
}

% Portada UTP
\setbeamertemplate{title page}{%
  \begin{tikzpicture}[remember picture,overlay]
    \node[anchor=center,inner sep=0pt] at (current page.center) {%
      \includegraphics[width=\paperwidth,height=\paperheight]{utp-cover-bg.jpeg}%
    };
    \node[anchor=north,align=center,text=utpNegro,inner sep=0pt,
      text width=145mm,font=\usebeamerfont{title}]
      at ([yshift=-25.2mm]current page.north) {\inserttitle};
    \node[anchor=north,align=center,text=utpNegro,inner sep=0pt,
      text width=142mm,font=\usebeamerfont{subtitle}]
      at ([yshift=-38.2mm]current page.north) {\insertsubtitle};
    \node[anchor=north,align=center,text=utpNegro,inner sep=0pt,
      font=\usebeamerfont{date}]
      at ([yshift=-50.4mm]current page.north) {\insertdate};
    \node[anchor=north,align=center,text=utpNegro,inner sep=0pt,
      font=\usebeamerfont{author}]
      at ([yshift=-58.4mm]current page.north) {\insertauthor};
  \end{tikzpicture}%
}

% ---------- Comandos de apoyo didáctico ----------
\newcommand{\concepto}[2]{%
  \begin{block}{#1}\large #2\end{block}}
\newcommand{\pregunta}[1]{%
  \begin{alertblock}{Decide y explica}#1\end{alertblock}}
\newcommand{\analogia}[1]{%
  \begin{exampleblock}{En palabras simples}\small #1\end{exampleblock}}
\newcommand{\enlarealidad}[1]{%
  \begin{block}{Cómo se ve en la realidad}\small #1\end{block}}
\newcommand{\ojo}[1]{%
  \begin{alertblock}{Ojo con esto}\small #1\end{alertblock}}
\newcommand{\apoyo}{{\normalfont\footnotesize\color{utpGris}~· material de apoyo}}
\newcommand{\marca}[1]{\tikz[baseline=(m.base)]{%
  \node[rounded corners=1.5pt,fill=utpGrisClaro,draw=utpGrisMedio,inner xsep=3pt,inner ysep=1.5pt,
        font=\footnotesize] (m) {#1};}}
\newenvironment{terminal}{%
  \begin{beamercolorbox}[wd=\linewidth,sep=5pt,rounded=true]{terminal}%
  \ttfamily\small\ignorespaces}{%
  \end{beamercolorbox}}

\tikzset{
  caja/.style={draw=utpNegro, very thick, rounded corners=1.5pt, fill=white,
    align=center, minimum height=0.72cm, inner sep=5pt, font=\small},
  cajaazul/.style={caja, draw=utpGris, fill=white},
  cajaverde/.style={caja, draw=utpNegro, fill=utpGrisClaro},
  cajanaranja/.style={caja, draw=utpRojo!75, fill=utpRojo!7},
  flecha/.style={-{Stealth[length=2.3mm]}, very thick, draw=utpNegro},
  flechavuelta/.style={{Stealth[length=2.3mm]}-{Stealth[length=2.3mm]}, very thick, draw=utpNegro}
}

\graphicspath{{../../Sesion05/slides/assets/}}



\graphicspath{{assets/}{../../Sesion05/slides/assets/}}
\usepackage{colortbl}
\usepackage{listings}
% Un manifiesto: las claves en rojo.
\lstset{basicstyle=\ttfamily\scriptsize,backgroundcolor=\color{utpGrisClaro},
  frame=none,aboveskip=4pt,belowskip=4pt,showstringspaces=false,
  columns=fullflexible,keepspaces=true,breaklines=true,
  keywords={kind,metadata,name,spec,replicas,selector,matchLabels,template,labels,containers,image},
  keywordstyle=\color{utpRojo}\bfseries,
  commentstyle=\color{utpGris}\itshape,morecomment=[l]{\#},
  xleftmargin=5pt,xrightmargin=5pt,framesep=5pt}
% Una terminal: fondo oscuro, sin colorear palabras.
\lstdefinestyle{terminal}{basicstyle=\ttfamily\scriptsize\color{white},backgroundcolor=\color{utpNegro},
  keywords={},morecomment=[l]{\#},commentstyle=\color{utpGrisMedio},upquote=true}
\usepackage{upquote}
\usepackage{pifont}
\definecolor{okVerde}{HTML}{2E7D32}
\usetikzlibrary{decorations.pathreplacing,positioning}

\title{Integración de Servicios Cloud GCP}
\subtitle{Redes, volúmenes y una introducción a Kubernetes}
\date{Semana 15 - Sesión 30}
\author{Rel Guzman}

% Íconos simples para los diagramas.
\tikzset{
 caja/.style={draw=utpRojo,thick,rounded corners=4pt,align=center},
 zona/.style={draw=utpGrisMedio,dashed,rounded corners=4pt},
 flecha/.style={-latex,thick},
 pod/.style={draw=utpRojo,thick,rounded corners=3pt,fill=utpRojo!30,align=center,font=\scriptsize\ttfamily,minimum width=1.5cm,minimum height=.6cm},
}

\begin{document}

{\setbeamertemplate{footline}{}\maketitle}

% 2
\begin{frame}{Revisión: lo que vimos en la Sesión 29}
\begin{columns}[T,onlytextwidth]
\column{.46\textwidth}
\begin{block}{Docker Compose}
\small Un archivo, \texttt{compose.yaml}, con los servicios \texttt{web}, \texttt{api} y \texttt{db}. \texttt{docker compose up -d} los levanta.
\end{block}
\begin{block}{Redes}
\small Cada servicio se alcanza por su nombre: \texttt{api:8080}, \texttt{db:5432}.
\end{block}
\begin{block}{Volúmenes}
\small Los pedidos quedan en el volumen \texttt{datos}, aparte del contenedor.
\end{block}
\column{.50\textwidth}
\begin{alertblock}{Responde}
\small
1. ¿Por qué la API busca la base en \texttt{db:5432} y no en \texttt{localhost:5432}?\par\medskip
2. ¿Qué servicio de Café Nube publica un puerto a tu computadora?\par\medskip
3. ¿Qué diferencia hay entre \texttt{docker compose down} y \texttt{docker compose down -v}?
\end{alertblock}
\end{columns}
\end{frame}

% 3
\begin{frame}{Revisión: la pregunta de hoy}
\begin{center}
\begin{tikzpicture}[font=\small]
  \node[caja,fill=utpGrisClaro,text width=1.6cm,minimum height=.9cm] (web) at (0,0) {\textbf{web}};
  \node[caja,fill=utpRojo!30,text width=1.6cm,minimum height=.9cm] (api) at (2.8,0) {\textbf{api}};
  \node[caja,draw=okVerde,fill=okVerde!12,text width=1.6cm,minimum height=.9cm] (db) at (5.6,0) {\textbf{db}};
  \node[zona,fit=(web)(db),inner sep=8pt,label={[font=\scriptsize,text=utpGris]below:una sola computadora}] (pc) {};
  \draw[flecha] (web) -- (api);
  \draw[flecha] (api) -- (db);
  \node[align=left,font=\small,text width=4.6cm,anchor=west] at (7.4,0) {¿Y si esa computadora se apaga?\\[2pt] ¿Y si un lunes llegan diez veces más pedidos?};
\end{tikzpicture}
\end{center}
\vspace{3mm}
\begin{columns}[T,onlytextwidth]
\column{.32\textwidth}
\begin{block}{1. Redes y volúmenes}
\small ¿Qué IP tiene cada contenedor? ¿Dónde quedan los archivos que copias?
\end{block}
\column{.32\textwidth}
\begin{block}{2. Clúster}
\small ¿Cómo se usan varias máquinas como si fueran una?
\end{block}
\column{.32\textwidth}
\begin{block}{3. Kubernetes}
\small ¿Quién vuelve a arrancar la API si se cae, sin esperar a una persona?
\end{block}
\end{columns}
\end{frame}

% 4
\begin{frame}[c]
\begin{block}{Logro de la sesión}
Al finalizar, podrás explicar con tus palabras cómo se encuentran los contenedores en una \textbf{red}, dónde quedan sus \textbf{archivos} y en qué se diferencian de un \textbf{volumen}, qué son un \textbf{clúster} y un \textbf{nodo}, y cómo \textbf{Kubernetes} corre la API de Café Nube en varias copias que se reparan solas: con \textbf{kind} en tu computadora y con \textbf{GKE} en Google Cloud.
\end{block}
\end{frame}

% 5
\begin{frame}{Saberes previos}
\begin{itemize}\setlength\itemsep{14pt}
  \item En el Dockerfile de la API, ¿qué instrucción mete \texttt{api.py} en la imagen?
  \item ¿Cuántas copias de la API corre \texttt{docker compose up -d}?
  \item Si la computadora donde corre Café Nube se apaga, ¿quién la vuelve a levantar?
\end{itemize}
\end{frame}

% 6
\begin{frame}[c]{Utilidad de esta sesión}
\begin{center}
\begin{minipage}{.84\textwidth}
{\large Con Compose, Café Nube corre en \textbf{una computadora} y en \textbf{una copia}. Las apps reales corren en \textbf{varias máquinas} y en \textbf{varias copias}, para que no se caigan cuando una falla. Eso hace \textbf{Kubernetes}, y en Google Cloud se usa con \textbf{GKE}.}
\end{minipage}
\end{center}
\end{frame}

% 7
\begin{frame}[fragile]{Redes: cada contenedor tiene su IP}
\begin{center}
\begin{tikzpicture}[font=\small]
  \node[caja,fill=utpGrisClaro,text width=2.1cm,minimum height=1.0cm] (web) at (0,0) {\textbf{web}\\\footnotesize\texttt{172.24.0.4}};
  \node[caja,fill=utpRojo!30,text width=2.1cm,minimum height=1.0cm] (api) at (4.2,0) {\textbf{api}\\\footnotesize\texttt{172.24.0.3}};
  \node[caja,draw=okVerde,fill=okVerde!12,text width=2.1cm,minimum height=1.0cm] (db) at (8.4,0) {\textbf{db}\\\footnotesize\texttt{172.24.0.2}};
  \node[zona,fit=(web)(db),inner sep=10pt,label={[font=\scriptsize,text=utpGris]below:red de Docker que crea Compose}] {};
  \draw[flecha] (web) -- node[above,font=\scriptsize\ttfamily]{api:8080} (api);
  \draw[flecha] (api) -- node[above,font=\scriptsize\ttfamily]{db:5432} (db);
\end{tikzpicture}
\end{center}
\vspace{1mm}
\begin{columns}[T,onlytextwidth]
\column{.44\textwidth}
\begin{block}{Dirección IP}
\small El número que identifica a un contenedor en la red. Docker \textbf{traduce el nombre} \texttt{db} a la IP de su contenedor.
\end{block}
\begin{alertblock}{Ojo}
\small La IP \textbf{cambia} si el contenedor se recrea. Por eso se usa el nombre, no la IP.
\end{alertblock}
\column{.52\textwidth}
\textbf{\small ¿Cómo veo la IP de un contenedor?}
\begin{lstlisting}[style=terminal]
docker inspect practica_cafe_nube-api-1
...
  "Networks": {
    "practica_cafe_nube_default": {
      ...
      "IPAddress": "172.24.0.3",
\end{lstlisting}
{\scriptsize Busca \texttt{IPAddress} en lo que responde. En la Sesión 29, los logs de la API decían \texttt{172.24.0.4}: era \texttt{web} pidiendo el menú.}
\end{columns}
\end{frame}

% 8
\begin{frame}{Volúmenes: ¿dónde quedan los archivos que copias?}
\begin{columns}[T,onlytextwidth]
\column{.46\textwidth}
\vspace{6mm}
\begin{center}
\begin{tikzpicture}[font=\scriptsize,capa/.style={draw=utpGrisMedio,fill=utpGrisClaro,minimum width=4.0cm,minimum height=.55cm,align=center}]
  \node[capa] (c1) at (0,0) {\texttt{FROM python:3.12-slim}};
  \node[capa,above=0pt of c1] (c2) {\texttt{RUN pip install ...}};
  \node[capa,above=0pt of c2] (c3) {\texttt{COPY . .}\quad (\texttt{api.py})};
  \node[capa,draw=utpRojo,fill=utpRojo!30,above=4pt of c3] (c4) {\textbf{capa del contenedor}\\lo que escribe al correr};
  \draw[decorate,decoration={brace,amplitude=4pt,mirror},thick] ([xshift=-3pt]c3.north west) -- node[left=6pt,align=right]{imagen\\solo lectura} ([xshift=-3pt]c1.south west);
\end{tikzpicture}
\end{center}
\column{.50\textwidth}
\begin{block}{\texttt{COPY} en el Dockerfile}
\small Copia al \textbf{construir}: \texttt{api.py} queda \textbf{dentro de la imagen}, en una capa. Todo contenedor de esa imagen lo tiene.
\end{block}
\begin{block}{\texttt{docker cp} a un contenedor que corre}
\small El archivo va a la \textbf{capa del contenedor}, encima de la imagen. \texttt{docker rm} la borra, con el archivo.
\end{block}
{\scriptsize Docker guarda todo en su carpeta, \texttt{/var/lib/docker}. En Windows y Mac está dentro de la máquina Linux que arranca Docker Desktop: no la ves en el Explorador.}
\end{columns}
\end{frame}

% 9
\begin{frame}[fragile]{Volúmenes: ¿\texttt{COPY} es lo mismo que un volumen? No}
\renewcommand{\arraystretch}{1.25}
\small
\begin{center}
\begin{tabular}{@{}>{\raggedright\arraybackslash}p{3.2cm}>{\raggedright\arraybackslash}p{4.6cm}>{\raggedright\arraybackslash}p{4.6cm}@{}}
\toprule
 & \textbf{\texttt{COPY} en el Dockerfile} & \textbf{Volumen}\\
\midrule
¿Cuándo se llena? & Al construir la imagen & Mientras corre el contenedor\\
¿Dónde queda? & Dentro de la imagen & Aparte, en su propia carpeta\\
¿Si borras el contenedor? & Sigue en la imagen; lo que cambió adentro se pierde & \textcolor{okVerde}{\textbf{Sigue}}\\
¿Para qué? & El código: \texttt{api.py}, \texttt{index.html} & Los datos: los pedidos de la base\\
\bottomrule
\end{tabular}
\end{center}
\vspace{1mm}
\textbf{\small ¿Dónde está el volumen?}
\begin{lstlisting}[style=terminal]
docker volume inspect practica_cafe_nube_datos
    "Mountpoint": "/var/lib/docker/volumes/practica_cafe_nube_datos/_data",
\end{lstlisting}
\end{frame}

% 10
\begin{frame}{Clúster y nodo: ¿de dónde vienen?}
\begin{columns}[T,onlytextwidth]
\column{.52\textwidth}
\begin{block}{Clúster}
\small Del inglés \textit{cluster}, «grupo» o «racimo». En informática, desde los años 80: \textbf{varias computadoras conectadas que trabajan como una sola}.
\end{block}
\begin{block}{Nodo}
\small Del latín \textit{nodus}, «nudo»: en el dibujo de una red, el punto donde llegan las líneas. En un clúster, \textbf{cada máquina}, física o virtual.
\end{block}
\begin{block}{Kubernetes (K8s)}
\small Del griego \textit{kybernētēs}, «timonel»: por eso su logo es un timón. Lo creó \textbf{Google} y lo liberó en 2014. K8s: la K, ocho letras y la s.
\end{block}
\column{.44\textwidth}
\begin{center}
\begin{tikzpicture}[font=\scriptsize]
  \node[caja,fill=utpGrisClaro,text width=3.6cm,minimum height=.8cm] (cp) at (0,1.9) {\textbf{plano de control}\\decide dónde va cada cosa};
  \foreach \i/\x in {1/-1.6,2/0,3/1.6} {
    \node[caja,text width=1.2cm,minimum height=1.3cm] (n\i) at (\x,0) {};
    \node[font=\scriptsize\bfseries,anchor=north] at (n\i.north) {nodo \i};
    \draw[flecha,utpGrisMedio] (cp.south) -- (n\i.north);
  }
  \node[pod,minimum width=.9cm,minimum height=.35cm] at (-1.6,-.2) {api};
  \node[pod,minimum width=.9cm,minimum height=.35cm] at (0,-.2) {api};
  \node[pod,minimum width=.9cm,minimum height=.35cm,fill=utpGrisClaro] at (1.6,-.2) {web};
  \node[zona,fit=(cp)(n1)(n3),inner sep=6pt,label={[font=\scriptsize,text=utpGris]below:clúster}] {};
\end{tikzpicture}
\end{center}
\end{columns}
\end{frame}

% 11
\begin{frame}{Kubernetes: ¿por qué?}
\begin{columns}[T,onlytextwidth]
\column{.48\textwidth}
\begin{block}{Con Compose, en una computadora}
\small
\begin{itemize}\setlength\itemsep{3pt}
  \item Se apaga la máquina: la app se cae y alguien la levanta a mano.
  \item Llegan diez veces más pedidos: una copia de la API no alcanza.
\end{itemize}
\end{block}
\column{.48\textwidth}
\begin{exampleblock}{Orquestador}
\small Un programa que \textbf{reparte contenedores entre varias máquinas y los vigila}: si uno se cae, lo vuelve a arrancar; si hacen falta más, crea más. El más usado es \textbf{Kubernetes}.
\end{exampleblock}
\end{columns}
\vspace{4mm}
\begin{center}
\begin{tikzpicture}[font=\small]
  \node[caja,text width=1.6cm,minimum height=1.2cm] (n1) at (0,0) {};
  \node[font=\scriptsize\bfseries,anchor=north] at (n1.north) {nodo 1};
  \node[pod,minimum width=1cm,minimum height=.35cm] at (0,-.2) {api};
  \draw[utpRojo,very thick] (n1.north west) -- (n1.south east);
  \draw[utpRojo,very thick] (n1.north east) -- (n1.south west);
  \node[caja,text width=1.6cm,minimum height=1.2cm] (n2) at (7,0) {};
  \node[font=\scriptsize\bfseries,anchor=north] at (n2.north) {nodo 2};
  \node[pod,minimum width=1cm,minimum height=.35cm] at (7,-.2) {api};
  \draw[flecha,dashed] (n1.east) -- node[above=2pt,font=\scriptsize]{Kubernetes la arranca en otro nodo} (n2.west);
\end{tikzpicture}
\end{center}
\end{frame}

% 12
\begin{frame}{Pods y manifiestos: describes lo que quieres}
\begin{center}
\begin{tikzpicture}[font=\small]
  \node[caja,dashed,text width=2.6cm,minimum height=1.0cm] (yo) at (0,0) {\textbf{Tú}\\\footnotesize\texttt{deployment.yaml}};
  \node[caja,fill=utpGrisClaro,text width=2.2cm,minimum height=1.0cm] (cp) at (5,0) {\textbf{plano de control}};
  \node[caja,text width=2.4cm,minimum height=1.3cm] (nodo) at (9.6,0) {};
  \node[font=\scriptsize\bfseries,anchor=north] at (nodo.north) {nodo};
  \node[pod] at (9.6,-.2) {pod: api};
  \draw[flecha] (yo) -- node[above,font=\scriptsize\ttfamily]{kubectl apply -f} (cp);
  \draw[flecha] (cp) -- node[above,font=\scriptsize]{crea} (nodo);
\end{tikzpicture}
\end{center}
\vspace{2mm}
\begin{columns}[T,onlytextwidth]
\column{.48\textwidth}
\begin{block}{Manifiesto}
\small Un archivo YAML que \textbf{describe} lo que quieres en el clúster. Con Docker dabas órdenes una por una; aquí describes cómo debe quedar todo, y Kubernetes se encarga.
\end{block}
\column{.48\textwidth}
\begin{block}{Pod}
\small La unidad más pequeña de Kubernetes: \textbf{uno o más contenedores que viven juntos}, en el mismo nodo y con la misma IP. Casi siempre, uno solo.
\end{block}
\end{columns}
\end{frame}

% 13
\begin{frame}[fragile]{Deployment: copias que se reparan solas}
\begin{columns}[T,onlytextwidth]
\column{.38\textwidth}
\begin{lstlisting}
kind: Deployment
metadata:
  name: api
spec:
  replicas: 3
  selector:
    matchLabels:
      app: api
  template:       # cada pod
    metadata:
      labels:
        app: api
    spec:
      containers:
        - name: api
          image: cafe-api:1
\end{lstlisting}
{\scriptsize Resumido: el archivo completo está en la guía.}
\column{.58\textwidth}
\begin{block}{Deployment}
\small Mantiene corriendo un número de \textbf{réplicas}, copias iguales de un pod. Si falta una, crea otra; si sobra, la borra.
\end{block}
\begin{block}{Label y selector}
\small \texttt{app: api} es una \textbf{label}, una marca en cada pod. El \texttt{selector} cuenta los pods que la llevan.
\end{block}
\begin{center}
\begin{tikzpicture}[font=\scriptsize]
  \node[pod,minimum width=1.1cm] (a) at (0,0) {api};
  \node[pod,minimum width=1.1cm] (b) at (1.4,0) {api};
  \node[pod,minimum width=1.1cm,fill=white,dashed] (c) at (2.8,0) {api};
  \draw[utpRojo,very thick] (c.north west) -- (c.south east);
  \draw[flecha] (3.5,0) -- (4.3,0);
  \node[pod,minimum width=1.1cm,draw=okVerde,fill=okVerde!12] at (5.0,0) {api};
  \node[font=\scriptsize,text=utpGris] at (2.5,-.6) {se borra uno: el Deployment ve 2 de 3 y crea otro};
\end{tikzpicture}
\end{center}
\end{columns}
\end{frame}

% 14
\begin{frame}{Service: un nombre fijo para pods que cambian}
\begin{center}
\begin{tikzpicture}[font=\small]
  \node[caja,fill=utpGrisClaro,text width=1.6cm,minimum height=.9cm] (web) at (0,0) {\textbf{pod web}};
  \node[caja,fill=white,text width=2.0cm,minimum height=.9cm] (svc) at (4.2,0) {\textbf{Service}\\\footnotesize\texttt{api}};
  \node[pod] (p1) at (8.4,.9) {10.244.0.5};
  \node[pod] (p2) at (8.4,0) {10.244.0.6};
  \node[pod] (p3) at (8.4,-.9) {10.244.0.7};
  \node[font=\scriptsize,text=utpGris,anchor=west] at (9.3,0) {pods con \texttt{app: api}};
  \draw[flecha] (web) -- node[above,font=\scriptsize\ttfamily]{api:8080} (svc);
  \draw[flecha] (svc) -- (p1);
  \draw[flecha] (svc) -- (p2);
  \draw[flecha] (svc) -- (p3);
\end{tikzpicture}
\end{center}
\vspace{1mm}
\begin{columns}[T,onlytextwidth]
\column{.48\textwidth}
\begin{block}{Service}
\small Cada pod nuevo trae \textbf{otra IP}. El Service tiene un \textbf{nombre fijo}, \texttt{api}, y reparte los pedidos entre los pods con su label. Los pods cambian; el Service queda.
\end{block}
\column{.48\textwidth}
\begin{alertblock}{De Compose a Kubernetes}
\small El servicio \texttt{api} de Compose se reparte en dos objetos: el \textbf{Deployment} corre las copias y el \textbf{Service} les da el nombre. Nginx sigue pidiendo a \texttt{api:8080}.
\end{alertblock}
\end{columns}
\end{frame}

% 15
\begin{frame}{kind y kubectl: Kubernetes en tu computadora}
\begin{columns}[T,onlytextwidth]
\column{.36\textwidth}
\begin{block}{kind}
\small \textit{Kubernetes IN Docker}: arma un clúster de práctica \textbf{dentro de Docker}. Cada nodo es un contenedor.
\end{block}
\begin{block}{kubectl}
\small El programa de terminal para hablarle a Kubernetes, como \texttt{docker} para Docker.
\end{block}
\begin{center}
\begin{tikzpicture}[font=\scriptsize]
  \node[caja,text width=3.2cm,minimum height=1.3cm] (nodo) at (0,0) {};
  \node[font=\scriptsize\ttfamily,anchor=north] at (nodo.north) {cafe-nube-control-plane};
  \node[pod,minimum width=.9cm,minimum height=.35cm] at (-.9,-.2) {api};
  \node[pod,minimum width=.9cm,minimum height=.35cm,fill=utpGrisClaro] at (.9,-.2) {web};
  \node[zona,fit=(nodo),inner sep=6pt,label={[font=\scriptsize,text=utpGris]below:Docker Desktop}] {};
\end{tikzpicture}
\end{center}
\column{.60\textwidth}
\renewcommand{\arraystretch}{1.15}
\begin{tabular}{@{}>{\ttfamily\scriptsize}l>{\raggedright\arraybackslash\scriptsize}p{2.3cm}@{}}
\toprule
\textsf{\textbf{\small Comando}} & \textbf{\small Qué hace}\\
\midrule
kind create cluster --name cafe-nube & Crea el clúster\\
kubectl get nodes & Lista los nodos\\
kubectl apply -f deployment.yaml & Entrega un manifiesto\\
kubectl get pods -o wide & Los pods, con su IP\\
kubectl logs NOMBRE-DEL-POD & Sus mensajes\\
kubectl delete pod NOMBRE-DEL-POD & Borra un pod\\
kubectl scale deployment api --replicas=5 & Pide cinco copias\\
kubectl port-forward service/web 8080:80 & Túnel, como \texttt{-p}\\
\bottomrule
\end{tabular}
\end{columns}
\end{frame}

% 16
\begin{frame}[fragile]{En Google Cloud: Cloud Build, Artifact Registry y GKE}
\begin{center}
\begin{tikzpicture}[font=\small]
  \node[caja,dashed,text width=2.0cm,minimum height=1.1cm] (cod) at (0,0) {\textbf{Tu código}\\\footnotesize\texttt{api.py}, Dockerfile};
  \node[caja,fill=utpGrisClaro,text width=2.3cm,minimum height=1.1cm] (cb) at (3.6,0) {\textbf{Cloud Build}\\\footnotesize construye la imagen};
  \node[caja,fill=utpRojo!30,text width=2.3cm,minimum height=1.1cm] (ar) at (7.4,0) {\textbf{Artifact Registry}\\\footnotesize guarda la imagen};
  \node[caja,draw=okVerde,fill=okVerde!12,text width=2.3cm,minimum height=1.1cm] (gke) at (11.2,0) {\textbf{GKE}\\\footnotesize la corre en pods};
  \draw[flecha] (cod) -- (cb);
  \draw[flecha] (cb) -- (ar);
  \draw[flecha] (ar) -- (gke);
\end{tikzpicture}
\end{center}
\vspace{2mm}
\begin{columns}[T,onlytextwidth]
\column{.50\textwidth}
\renewcommand{\arraystretch}{1.25}
\small
\begin{tabular}{@{}>{\raggedright\arraybackslash}p{2.9cm}>{\raggedright\arraybackslash}p{3.3cm}@{}}
\toprule
\textbf{En tu computadora} & \textbf{En Google Cloud}\\
\midrule
\texttt{docker build} & \textbf{Cloud Build}: construye en la nube\\
\texttt{kind load docker-image} & \textbf{Artifact Registry}: el registro de tus imágenes\\
El clúster de kind & \textbf{GKE}: Kubernetes administrado por Google\\
\bottomrule
\end{tabular}
\column{.46\textwidth}
\begin{block}{GKE}
\small Google maneja el plano de control; los nodos son máquinas virtuales de Compute Engine. Los \textbf{mismos manifiestos} sirven: solo cambia de dónde sale la imagen.
\end{block}
{\scriptsize Una imagen en Artifact Registry:\\\texttt{us-central1-docker.pkg.dev/}\\\texttt{\ \ mi-proyecto/cafe/cafe-api:1}}
\end{columns}
\end{frame}

% 17
\begin{frame}{Cierre: preguntas para comprobar}
\begin{block}{Responde con tus palabras}
\small
1. Borras el contenedor de la API y creas otro con la misma imagen. ¿Sigue \texttt{api.py}? ¿Y un archivo que le pasaste con \texttt{docker cp}?\\\medskip
2. Un Deployment pide \texttt{replicas: 3} y borras uno de sus pods. ¿Qué pasa?\\\medskip
3. ¿Por qué Nginx le pide a la API en \texttt{api:8080} y no en la IP de un pod?
\end{block}
\end{frame}

% 18
\begin{frame}{Referencias}
\scriptsize
\begin{itemize}\setlength\itemsep{4pt}
  \item Ruta \textit{Despliegue con Docker}, curso 2, \textit{Varios contenedores con Compose}: redes entre contenedores y volúmenes.
  \item Ruta \textit{Despliegue con Docker}, curso 3, \textit{Kubernetes desde cero}: clúster y nodos, pods, Deployments, Services y el escenario Café Nube en Kubernetes.
  \item Docker: \texttt{docs.docker.com}. Kubernetes: \texttt{kubernetes.io}. kind: \texttt{kind.sigs.k8s.io}.
  \item Google Cloud: \texttt{cloud.google.com/kubernetes-engine}, \texttt{cloud.google.com/artifact-registry} y \texttt{cloud.google.com/build}.
  \item Material previo del curso: Sesión 29 (Docker Compose, redes y volúmenes).
  \item Las respuestas de las preguntas están en \texttt{../respuestas.md}.
\end{itemize}
\end{frame}

\end{document}
